% 文件开头添加条件判断
\newif\ifstandalone
\standalonetrue  % 默认独立编译模式

% 检测是否在主文档中
\ifdefined\maindoc
  \standalonefalse
\fi

\ifstandalone
  \documentclass[UTF8,a4paper]{ctexart}
  \usepackage{amsmath,amssymb,amsthm,graphicx,hyperref}
  \usepackage{geometry}
  \geometry{a4paper,margin=2.5cm}
  \title{第1章　商用机器人操作器}
  \author{《机器人操作器控制：理论与实践》中文译本}
  \date{}
\fi

\ifstandalone
  \begin{document}
  \maketitle
\fi

\section{引言}

本章通过概述商用机器人操作器、传感器和控制器，为本书奠定基础。我们要强调的是，如果一个人希望获得高性能的柔性机器人工作单元，那么就有必要为机器人操作器设计如本书所述的先进控制系统。

在研究机器人控制、规划、传感器和人机接口的高级技术时，了解现有商用系统是很重要的。这允许人们在现有技术的背景下开发新技术，从而能够在现有机器人系统上实现这些新技术。

全国制造商协会的一份报告[NAM 1998]指出，20世纪90年代美国商业制造业成功的两个最重要驱动因素是可重构制造工作单元和工厂中的局域网。在本章中，我们将讨论柔性机器人工作单元、商用机器人构型、商用机器人控制器、到互联网的信息集成以及机器人工作单元传感器。有关这些主题的更多信息可以在《机械工程手册》[Lewis 1998]和《计算机科学工程手册》[Lewis and Fitzgerald 1997]中找到。

\subsection{柔性机器人工作单元}

在工厂自动化和其他领域，曾经普遍使用围绕传送带或其他运输系统建造的固定布局，其中每个机器人执行特定任务。这些装配线具有不同的工作站，每个工作站执行专用功能。机器人已被用于工作站级别执行装配、钻孔、表面处理、焊接、码垛等操作。在装配线中，零件由运输系统顺序地路由到工作站。此类系统安装成本非常高，需要一组工程专家来设计和编程，并且随着需求的变化极难修改或重新编程。在今天的高混合低批量（HMLV）制造场景中，这些特性宣告了这种刚性过时设计的死亡。

在装配线中，机器人被限制在放入刚性顺序系统中。机器人是具有多种功能的通用机器，通过将它们用作柔性机器人工作单元的基础，其潜力可以显著增加，如图1.1.1所示的UTA自动化与机器人测试单元。在柔性机器人工作单元中，机器人用于零件搬运、装配和其他加工操作。通过重新编程机器人，可以改变工作单元的整个功能。工作单元设计旨在充分利用机器人的工作空间，并且铣床、钻床、振动零件送料机等组件被放置在机器人的工作空间内，以便由机器人提供服务。与装配线相反，物理布局并不预先确定固定的操作或作业顺序。因此，随着产品需求的变化，只需要在软件中重新编程工作单元即可。工作单元非常适合制造业和其他领域中新兴的HMLV条件。

机器人工作单元的日益普及已经将重点从硬件设计转移开，并将新的重点放在包括规划、协调和控制（PC\&C）功能的创新软件技术和架构上。需要对机器人控制器进行大量研究，以使机器人具备现代柔性工作单元所需的灵活性、精度和功能。本书的其余部分详细介绍了此类先进控制技术。

\section{商用机器人构型与类型}

本节中的大部分信息由Mick Fitzgerald准备，他当时是UTA自动化与机器人研究所（ARRI）的经理。

机器人是高度可靠、可信和技术先进的工厂设备。世界上大多数机器人由使用可靠的现成组件技术的成熟公司提供。所有商用工业机器人都有两个物理上分离的基本元素——操作器手臂和控制器。大多数商用机器人的基本架构基本相同，由串行连杆运动机器上的数字伺服控制电动电机驱动组成，通常不超过六个轴（自由度）。所有机器人都配有专有控制器。几乎所有的机器人应用都需要工程师和技术人员进行大量的设计和实施工作。使每个机器人独特的是如何将组件组合在一起以实现性能，从而产生具有竞争力的产品。工业机器人应用中最重要的考虑集中在两个问题上：操作和集成。

\subsection{操作器性能}

运动学结构、轴驱动机构设计和实时运动控制的综合效果决定了主要操作性能特征：可达性和灵巧性、负载能力、快速性和精度。在根据制造商公布的性能规格做出决策和进行比较时必须谨慎，因为测量和报告这些规格的方法在行业内并不统一。通常使用运动测试、仿真或其他分析技术来验证每个应用的性能。

\textbf{可达性}通过测量机器人运动所描述的工作空间范围来表征，而\textbf{灵巧性}通过各个关节的角位移来表征。一些机器人在其可达性边界内会有无法使用的空间，如死区、奇异姿态和腕部缠绕姿态。

\textbf{负载重量}由所有工业机器人的制造商指定。一些制造商还指定了旋转腕轴的惯性负载。通常给出的负载是在极端速度和可达条件下。所有工具、工件、电缆和软管的重量和惯性都必须作为负载的一部分包括在内。

\textbf{快速性}在确定吞吐量方面至关重要，但很难从公布的机器人规格中确定。大多数制造商会指定单个关节或特定运动学工具点的最大速度。然而，工作周期中的平均速度是感兴趣的快速性特征。

\textbf{精度}通常通过测量重复性来表征。几乎所有机器人都指定了静态位置重复性。很少指定准确度，但它可能比重复性大四倍。动态精度，或在连续路径上跟踪位置、速度和加速度的重复性和准确度，通常不被指定。

\subsection{常见运动学构型}

所有常见的商用工业机器人都是串行连杆操作器，通常不超过六个运动学耦合轴。按照惯例，运动轴按从基座到腕部的顺序编号。前三个轴负责机器人的空间定位运动；它们的构型决定了机器人可以定位的空间形状。运动链中的任何后续轴通常提供旋转运动以定向机器人手臂的末端，被称为腕部轴。在机器人腕部中，三个轴通常相交，以三维定向的形式产生真正的独立定位。附录A提供了球形机器人腕部机构的运动学分析。请注意，在我们的三维空间中，完全独立的空间定位需要三个自由度，完全独立的定向定位也需要三个自由度。

机器人轴可以在其驱动连杆中产生两种主要类型的运动——旋转或移动。旋转关节是拟人化的（例如像人类关节），而移动关节能够像汽车收音机天线一样伸缩。通常根据其前三个轴的方向和类型对机器人进行分类是非常有用的。有四种非常常见的商用机器人构型：\textbf{关节型}、\textbf{I型SCARA}、\textbf{II型SCARA}和\textbf{笛卡尔型}。另外两种构型，\textbf{圆柱型}和\textbf{球形}，现在已不那么常见。

附录C包含了一些常见机器人操作器的动力学，用于本书的控制仿真。

\textbf{关节型手臂}。商用关节型手臂的种类繁多，大多数有六个轴（图1.2.1）。所有这些机器人的轴都是旋转的。第二和第三轴共面并协同工作以在垂直平面内产生运动。基座中的第一轴是垂直的并旋转手臂以扫出大的工作体积。已经设计了各种类型的驱动机构，以允许腕部和前臂驱动电机和齿轮箱安装在靠近第一和第二旋转轴的位置，从而最小化手臂的延伸质量。当需要五个或更多自由度时，设计良好的关节型手臂的工作空间效率（即相对于手臂大小的快速灵巧可达程度）是其他手臂构型无法比拟的。关节型手臂性能的一个主要限制因素是第二轴必须工作以提升随后的手臂结构和负载。从历史上看，关节型手臂未能达到与其他手臂构型一样高的精度，因为所有轴都有关节角度位置误差，这些误差乘以连杆半径并在整个手臂上累积。

\textbf{I型SCARA}。I型SCARA（选择性柔顺装配机器人手臂）手臂使用两个平行旋转关节在水平平面内产生运动。手臂结构承载重量，但第一和第二轴不做提升。I型SCARA的第三轴通过添加垂直或z轴提供工作体积。第四旋转轴将添加绕z轴的旋转以控制水平平面内的方向。这种类型的机器人很少发现超过四个轴。I型SCARA广泛用于电子元件和设备的装配，并广泛用于小型和中型机械装配的装配。

\textbf{II型SCARA}。II型SCARA也是一种四轴构型，与I型的不同之处在于第一轴是一个长的垂直移动z行程，提升两个平行旋转轴及其连杆。对于在更长距离（超过约三英尺）上快速移动更重负载（超过约75磅），II型SCARA构型比I型更有效。

\begin{quote}
\textit{注：读者可以通过}\href{https://frankjie09.github.io/Robot-Manipulator-Control-CN/figures/SCARA/SCARA_Type_I_vs_II_Comparison.html}{\textbf{交互式可视化工具}}\textit{对比I型和II型SCARA机器人的运动学差异。}
\end{quote}

\textbf{笛卡尔坐标机器人}。笛卡尔坐标机器人使用正交移动轴，通常称为x、y和z，通过其矩形工作空间平移其末端执行器或负载。可以添加一个、两个或三个旋转腕轴用于定向。商业机器人公司提供几种类型的笛卡尔坐标机器人，工作空间大小从几立方英寸到数万立方英尺，负载能力可达数百磅。门式机器人具有高架桥结构，是最常见的笛卡尔型，非常适合需要在大面积和/或大负载上提供服务的物料搬运应用。它们在弧焊、水射流切割和大型复杂精密零件的检查等应用中特别有用。

模块化笛卡尔机器人也可从几家商业来源获得。每个模块是一个自包含的完全功能的单轴执行器；这些模块可以定制组装用于特殊用途应用。

\textbf{球型和圆柱坐标机器人}。球型坐标机器人的前两个轴是旋转的且相互正交，第三轴提供移动径向延伸。结果是具有球型工作体积的自然球坐标系统。圆柱坐标机器人的第一轴是旋转基座旋转。第二和第三是移动的，产生自然的圆柱运动。球型和圆柱型机器人的商业型号最初非常常见，在机器看护和物料搬运应用中很受欢迎。数百台仍在使用，但现在只有少数商用型号。这两种构型使用下降的原因归因于使用移动连杆进行径向延伸/缩回运动的问题；实心臂需要间隙才能完全缩回。

\textbf{并联连杆操作器}。对于某些特殊用途应用，并联连杆机器人比串联连杆机器人更适合。这些机器人通常有三个或六个并联连杆，每个连杆连接到一个固定基座和一个移动工作平台。如图1.2.7所示。通过适当的设计，六连杆并联连杆操作器可以使工作平台具有六自由度运动。军用Link训练器是一种大型并联连杆机器人，移动飞行员座椅。这些机器人比串联连杆机器人具有更大的刚度和精度，在串联连杆机器人中，每个连杆的定位误差随着从基座向外移动而累积。因此，轻质并联连杆机器人能够精确地移动大负载。这些机器人已被用于加工和表面精加工精密工业和航空航天部件，如舱壁和飞行器外壳。

并联连杆机器人是一个闭链运动学系统，因此相对难以分析[Liu and Lewis 1993]。这些机器人的控制系统设计问题更加困难。

\begin{quote}
\textit{注：读者可以通过}\href{https://frankjie09.github.io/Robot-Manipulator-Control-CN/figures/Parallel_Link/Parallel_Link_Manipulator.html}{\textbf{交互式可视化工具}}\textit{探索六自由度并联连杆操作器（Stewart平台）的运动特性。}
\end{quote}

\subsection{商用机器人的驱动类型}

绝大多数商用工业机器人使用带减速传动的电动伺服电机驱动。交流和直流电机都很流行。现在有一些用于喷漆应用的伺服液压关节型手臂机器人。很少发现带伺服气动驱动轴的机器人。使用所有类型的机械传动，但趋势是向低间隙和零间隙型驱动发展。一些机器人使用直接驱动方法来消除与其他驱动相关的惯性和机械间隙放大。关节角度位置传感器是实时伺服级控制所需的，通常被认为是传动系统的重要部分。较少提供速度反馈传感器。

\begin{quote}
\textit{注：读者可以通过}\href{https://frankjie09.github.io/Robot-Manipulator-Control-CN/figures/Gear_Drive_Comparison.html}{\textbf{交互式可视化工具}}\textit{对比传统齿轮传动与直接驱动方式的差异，观察间隙对控制精度的影响。}
\end{quote}

\section{商用机器人控制器}

商用机器人控制器是专门的多处理器计算系统，提供四个基本过程，允许将机器人集成到自动化系统中：运动轨迹生成和跟踪、运动/过程集成和排序、人机集成以及信息集成。

\subsection{运动轨迹生成和跟踪}

工业机器人运动生成有两个重要的与控制器相关的方面。一个是可编程的操作程度，另一个是执行受控编程运动的能力。每个机器人系统的独特方面其实时伺服级运动控制。实时控制的细节通常不会向用户透露，出于安全和专有信息保密的原因。每个机器人控制器通过其操作系统程序，将来自更高级协调器的数字数据通过精确计算和高速分配和通信转换为协调的手臂运动，各个轴运动命令由各个关节伺服控制器执行。大多数商用机器人控制器以16毫秒的采样周期运行。实时运动控制器 invariably 使用经典独立关节比例-积分-微分（PID）控制或PID的简单改进。这使得商用控制器适用于点对点运动，但大多数不适用于跟踪连续位置/速度轨迹或施加规定的力，即使有大量的编程工作。

最近，出现了更先进的控制器。Adept Windows系列的自动化控制器集成了机器人学、运动控制、机器视觉、力传感和制造逻辑在一个与Windows 98 \& Windows NT/2000兼容的单一控制平台上。Adept运动控制器可以配置为控制其他机器人和定制机构，并且是各种OEM系统的标准配置。

\subsection{运动/过程集成和排序}

运动/过程集成涉及将操作器运动与过程传感器或其他过程控制器设备协调。最基本的过程集成是通过离散数字输入/输出（I/O）。例如，机器人控制器外部的机器控制器可能会发送一位信号，表示它已准备好被机器人装载。机器人控制器必须能够读取数字信号并使用该信号执行逻辑操作（if then、wait until、do until等）。也就是说，一些机器人控制器内置了一些可编程逻辑控制器（PLC）功能。与传感器（例如视觉）的协调也通常提供。

\subsection{人机集成}

控制器的人机接口对于机器人系统的快速设置和编程至关重要。大多数机器人控制器有两种类型的人机接口可用：计算机式CRT/键盘终端用于离线和编辑程序代码，以及示教器，这是便携式手动输入终端，用于通过触摸键或操纵杆以遥操作方式命令运动。示教器通常是定位机器人最有效的方式，控制器中的存储器可以在回放模式下播放示教位置以执行运动轨迹。通过实践，人类操作员可以快速示教一系列在回放模式下链接在一起的点。目前，大多数机器人应用依赖于编程阶段人类专业知识的集成，以成功规划和协调机器人运动。这些接口机制在无阻碍的工作空间中有效，在编程和执行之间没有变化。它们不允许在执行期间人机接口或适应变化的环境。

更 recent 的高级机器人接口技术基于行为编程，其中各种特定行为以低级别编程到机器人控制器中（例如拾取零件、插入机器卡盘）。然后由更高级别的机器主管根据人类操作员的规定对这些行为进行排序并指定其特定运动参数。这种方法在[Mireles and Lewis 2001]中使用。

\subsection{信息集成}

随着向增加灵活性和敏捷性的趋势影响机器人学，信息集成变得越来越重要。许多商用机器人控制器现在通过通信端口（例如RS-232）采用集成PC接口支持信息集成功能，或者通过直接连接到机器人控制器数据总线。最近的集成努力使得可以将机器人工作单元连接到互联网，以允许远程站点监控和控制。有许多技术可以实现这一点，其中最方便的是LabVIEW 6.1，它不需要用Java编程。

\section{传感器}

本节中的大部分信息由Kok-Meng Lee准备[Lewis 1998]。传感器和执行器[Tzou and Fukuda 1992]作为换能器发挥作用，高级工作单元规划、协调和控制系统通过它与组成工作单元的硬件组件接口。传感器是一个至关重要的元素，因为它们将物理设备的状态转换为适合输入到工作单元PC\&C控制系统的信号；不合适的传感器可能引入错误，无论PC\&C系统多么复杂或昂贵，都无法使正确操作成为可能，而传感器的创新选择可以使控制和协调问题变得更加容易。

传感器有许多不同类型和许多不同的用途。考虑到与生物系统的类比，\textbf{本体感受器}是设备内部的传感器，产生关于该设备内部状态的信息（例如机器人手臂关节角度传感器）。\textbf{外感受器}产生关于设备外部其他硬件的信息。传感器产生模拟或数字输出；数字传感器通常提供关于机器或资源状态的信息（夹持器打开或关闭、机器已装载、作业完成）。传感器产生PC\&C层次结构所有级别所需的输出，包括用于：
\begin{itemize}
    \item 伺服级反馈控制（通常是模拟本体感受器）
    \item 过程监控和协调（通常是数字外感受器或零件检查传感器如视觉）
    \item 故障和安全监控（通常是数字——例如接触传感器、气动压力损失传感器）
    \item 质量控制检查（通常是视觉或扫描激光）。
\end{itemize}

传感器输出数据通常必须进行处理，以将其转换为对PC\&C目的有意义的形式。传感器加所需的信号处理显示为虚拟传感器。它作为数据抽象发挥作用——一组数据加上对该数据的操作（例如相机、加帧采集卡、加信号处理算法如图像增强、边缘检测、分割等）。一些传感器，包括伺服级反馈控制所需的本体感受器，是其主机设备的组成部分，因此传感器数据的处理和数据的使用发生在该设备内；然后，传感器数据在伺服控制级别或机器协调级别合并。其他传感器，通常是视觉系统，在复杂性上与机器人操作器相当，并由作业协调器协调，作业协调器将其视为有价值的共享资源，其使用由某种优先级分配（例如分派）方案分配给需要它们的作业。所谓的主动感知提出了一个有趣的协调问题，例如，机器人可能持有一个扫描相机，相机有效地接管运动协调问题，指挥机器人向何处移动以在后续图像中实现熵的最大减少（信息增加）。

\subsection{传感器类型}

本节从操作角度概述传感器。有关功能和物理原理的更多信息可以在[Fraden 1993]、[Fu et al. 1987]、[Snyder 1985]中找到。

\textbf{触觉传感器}。触觉传感器依靠与外部物体的物理接触。数字传感器如限位开关、微动开关和真空设备提供关于是否发生接触的二进制信息。有传感器可以检测滑移的开始。模拟传感器如弹簧加载杆提供更多信息。基于橡胶状碳或硅基弹性体并嵌入电气或机械组件的触觉传感器可以提供关于零件几何形状、位置等的非常详细的信息。弹性体可以包含电阻或电容元件，其电气特性随着弹性体压缩而变化。基于LSI技术的设计可以生产触觉网格垫，例如在单个垫上有64×64的"力点"。此类传感器产生具有类似相机数字图像属性的"触觉图像"，需要类似的数据处理。额外的触觉传感器属于随后讨论的"力传感器"分类。

\textbf{接近和距离传感器}。非接触接近传感器包括基于霍尔效应的设备或基于电磁效应的感应设备，可以检测约5毫米范围内的铁磁材料。此类传感器通常是数字的，产生关于物体是否在附近的二进制信息。基于电容的传感器检测约5毫米范围内的任何附近固体或液体。光学和超声传感器具有更长的范围。

距离传感器包括飞行时间测距设备如声纳和激光。商用Polaroid声纳提供约1英寸的精度，可达5英尺，角扇区精度约15度。对于移动机器人导航应用中的360度覆盖，可以使用扫描声纳和环形安装的多声纳。声纳通常有噪声，有虚假读数，需要低通滤波和其他旨在降低虚警率的数据处理。更昂贵的激光测距仪在距离上极其精确，具有很高的角分辨率。

\textbf{位置、速度和加速度传感器}。线性位置测量设备包括线性电位计和刚刚讨论的声纳和激光测距仪。线性速度传感器可以是基于激光或声纳的多普勒效应设备。

关节角度位置和速度本体感受器是机器人手臂伺服控制驱动轴的重要组成部分。角度位置传感器包括使用直流电压的电位计和使用交流电压的旋转变压器，精度为15分钟。光学编码器可以使用数字技术提供极高的精度。增量光学编码器使用三个光学传感器和单个交替不透明/透明区域环，如图1.4.1(a)所示，以提供相对于参考点的角位置信息和角速度信息；商用设备每转可能有1200个槽。更昂贵的绝对光学编码器，如图1.4.1(b)所示，有$n$个同心环的交替不透明/透明区域，需要$n$个光学传感器。如果采用格雷码，它们提供更高的精度并最小化与数据读取和传输相关的误差，其中连续扇区之间只有一位变化。精度为$360^\circ/2^n$，商用设备具有$n\approx 12$。

陀螺仪如果补偿与漂移相关的重复性问题，具有良好的精度。定向陀螺的精度约为1.5度。垂直陀螺的精度为0.5度，可用于测量多轴运动（例如俯仰和横滚）。速率陀螺直接测量速度，阈值约为0.05度/秒。

有各种基于应变计（下一段）、陀螺仪或晶体特性的加速度计可用。商用设备可用于测量沿三个轴的加速度。一种流行的新技术涉及微机电系统（MEMS），它要么是表面要么是体微加工设备。MEMS加速度计非常小、便宜、鲁棒且精确。MEMS传感器特别已用于汽车行业[Eddy 1998]。

\textbf{力和扭矩传感器}。有各种扭矩传感器可用，尽管通常不需要；例如，机器人手臂关节的内部扭矩可以从电机电枢电流计算。钻削工具上的扭矩传感器可以指示工具何时变钝。

线性力可以使用负载传感器或应变计测量。应变计是一种弹性传感器，其电阻是施加应变或变形的函数。压电效应，即在施加力时产生电压，也可用于力传感。其他力传感技术基于真空二极管、石英晶体（其共振频率随施加力变化）等。

机器人手臂力-扭矩腕部传感器在灵巧操作任务中非常有用。商用设备可以沿三个垂直轴测量力和扭矩，提供关于与腕部接触的笛卡尔力矢量$F$的完整信息。标准变换允许计算其他坐标中的力和扭矩。六轴力-扭矩传感器相当昂贵。

\textbf{光电传感器}。有各种光电传感器可用，有些基于光纤原理。这些响应速度约为50微秒，范围可达约45毫米，可用于检测零件和标签、扫描光学条形码、确认分拣任务中的零件通过等。

\textbf{其他传感器}。有各种传感器可用于测量压力、温度、流体流量等。这些在焊接等某些过程的闭环伺服控制应用中以及作业协调和/或安全中断程序中非常有用。

\subsection{传感器数据处理}

在任何传感器用于机器人工作单元之前，必须对其进行校准。根据传感器的不同，这可能涉及大量的实验、计算和安装后的调整工作。制造商通常提供校准程序，尽管在某些情况下，包括视觉，此类程序可能不明显，需要参考已发表的科学文献。系统修改后可能需要耗时的重新校准。

特别是对于更复杂的传感器如光学编码器，需要大量的传感器信号调理和处理。这可能包括信号放大、噪声抑制、数据从模拟到数字或从数字到模拟的转换等。制造商通常为这些目的提供硬件，应被视为机器人工作单元设计的传感器包的一部分。传感器及其信号处理硬件和软件算法可被视为数据抽象，被称为"虚拟传感器"。

如果需要解决信号处理问题，使用有限状态机（FSM）设计通常非常有用。来自增量光学编码器的典型信号如图1.4.2(a)所示；用于将其解码为角位置的FSM如图1.4.2(b)所示。FSM非常容易直接转换为逻辑门形式的硬件。用于声纳排序的FSM如图1.4.3(a)所示；从该FSM派生的声纳驱动硬件如图1.4.3(b)所示。

一个特定问题是从角位置测量获得角速度。很多时候，位置测量只是简单地差分，使用小采样周期来计算速度。如果信号中有任何噪声，这肯定会导致问题。几乎总是需要使用低通滤波导数，其中速度样本$v_k$从位置测量样本$p_k$使用以下公式计算：
\[
v_k = \frac{p_k - p_{k-1}}{T} + \alpha (v_{k-1} - v_k)
\]
其中$T$是采样周期，$\alpha$是小滤波系数。需要类似的方法来计算加速度。

\subsection{视觉系统、相机和照明}

典型的商用视觉系统符合20世纪50年代的RS-170标准，因此帧通过帧采集卡以每秒30帧的速率采集。图像是扫描的；在一个流行的美国标准中，每个完整的扫描或帧由525条线组成，其中480条包含图像信息。这个采样率和这种数量级的图像分辨率足以满足大多数应用，除了基于视觉的机器人手臂伺服。机器人视觉系统相机通常是电视相机——固态电荷耦合器件（CCD），对低于350纳米（紫外线）到1100纳米（近红外）的光波长响应，峰值响应约为800纳米，或电荷注入器件（CID），提供类似的光谱响应，峰值响应约为650纳米。线扫描CCD相机和面积扫描CCD相机都可用，分辨率范围在256到2048元素之间。中分辨率面积扫描相机产生256×256的图像，尽管现在已有1024×1024的高分辨率设备。线扫描相机适用于零件在相机前移动的应用，例如在传送带上。帧采集器通常支持多个相机，常见数量为四个，可支持黑白或彩色图像。

如果听任自然，机器人工作单元的照明可能会导致操作中的严重问题。常见问题包括低对比度图像、镜面反射、阴影和无关细节。这些问题可以通过过于复杂的图像处理来纠正，但所有这些问题都可以通过在工作单元设计阶段对细节进行适当关注来避免。照明技术包括光谱滤波、选择合适的照明源光谱特性、漫射照明技术、背光（产生易于处理的轮廓）、结构光（提供额外的深度信息并简化物体检测和解释）以及定向照明。

\section*{参考文献}

\begin{enumerate}
    \item[{[Decelle 1988]}] Decelle, L.S., "Design of a Robotic Workstation For Component Insertions," AT\&T Technical Journal, March/April 1988, Volume 67, Issue 2. pp 15–22.
    \item[{[Eddy 1998]}] Eddy, D.S., and D.R.Sparks, "Application of MEMS technology in automotive sensors and actuators," Proc. IEEE, vol. 86, no. 8, pp. 1747–1755, Aug. 1998.
    \item[{[Fraden 1993]}] Fraden, J. AIP Handbook Of Modern Sensors, Physics, Design, and Applications, American Institute of Physics. 1993.
    \item[{[Fu et al. 1987]}] Fu, K.S., R.C.Gonzalez, and C.S.G.Lee, \textit{Robotics}, McGraw-Hill, New York, 1987.
    \item[{[Jamshidi et al. 1992]}] Jamshidi, M., Lumia, R., Mullins, J., and Shahinpoor, M., 1992. \textit{Robotics and Manufacturing: Recent Trends in Research, Education, and Applications}, Vol. 4, ASME Press, New York.
    \item[{[Lewis and Fitzgerald 1997]}] Lewis, F.L., M.Fitzgerald, and K.Liu "Robotics," in \textit{The Computer Science and Engineering Handbook}, Allen B.Tucker, Jr. ed., Chapter 33, CRC Press, 1997.
    \item[{[Lewis 1998]}] Lewis, F.L., "Robotics," in \textit{Handbook of Mechanical Engineering}, F.Kreith ed., chapter 14, CRC Press, 1998.
    \item[{[Liu and Lewis 1993]}] Liu, K., F.L.Lewis, G.Lebret, and D.Taylor, "The singularities and dynamics of a Stewart Platform Manipulator," J. Intelligent and Robotic Systems, vol. 8, pp. 287–308, 1993.
    \item[{[Mireles and Lewis 2001]}] Mireles, J., and F.L.Lewis, "Intelligent Material Handling: Development and implementation of a matrix-based discrete-event controller," IEEE Trans. Industrial Electronics, vol. 48, no. 6, pp. 1087–1097, Dec. 2001.
    \item[{[NAM 1998]}] National Assoc. Manufacturers Report, "Technology on the Factory Floor III," ed. P.M.Swadimass, NAM Pub. Center, 1–800–637–3005, Aug. 1998.
    \item[{[Pugh 1983]}] Pugh, A., ed., 1983. \textit{Robotic Technology}, IEE Control Engineering Series 23, Pergrinus, London.
    \item[{[Snyder 1985]}] Snyder, W.E., 1985, \textit{Industrial Robots: Computer Interfacing and Control}, Prentice-Hall, Inc., Englewood Cliffs, New Jersey, USA.
    \item[{[Tzou and Fukuda 1992]}] Tzou, H.S., and Fukuda, T., \textit{Precision Sensors, Actuators, and Systems}, Kluwer Academic, 1992.
\end{enumerate}

\ifstandalone
  \end{document}
\fi
